\documentclass[twocolumn,10pt]{article}

\usepackage[margin=0.75in]{geometry}
\usepackage{amsmath,amssymb,amsthm}
\usepackage{graphicx}
\usepackage{xcolor}
\usepackage[colorlinks=true,linkcolor=blue,citecolor=blue,urlcolor=blue]{hyperref}
\usepackage{authblk}
\usepackage{abstract}

\renewcommand{\Affilfont}{\small}
\setlength{\affilsep}{0.3em}

\newtheorem{theorem}{Theorem}
\newtheorem{lemma}{Lemma}
\newtheorem{proposition}{Proposition}
\newtheorem{definition}{Definition}
\newcommand{\keff}{k_{\rm eff}}
\newcommand{\dDE}{\delta_{\rm DE}}
\newcommand{\dm}{\delta_{\rm m}}

\title{\vspace{-2em}\textbf{Which dark energy can structure growth see?
A response-kernel classification of geometric dark energy}}

\author[]{Eric Schultz}
\affil[]{Independent researcher \quad ORCID: 0009-0006-6283-1696}
\date{\today}

\begin{document}

\twocolumn[
  \begin{@twocolumnfalse}
  \maketitle
  \begin{abstract}
  \noindent
  A dark-energy model whose density is fixed by a geometric construction---a
  horizon length, a curvature scalar, a nonlocal functional of the
  metric---has no independent perturbation sector: its density contrast is
  \emph{inherited} from the gravitational potential through whatever functional
  defines it. We show that the observability of such a model in the linear
  structure-growth rate $f\sigma_8$ is controlled by a single quantity: the
  ultraviolet scaling of the kernel through which the dark-energy density
  responds to the potential. Because the density perturbation is always a linear
  functional of the metric perturbation, geometric dark-energy models partition
  exhaustively by this kernel's large-$k$ behaviour into a small number of
  classes. Local constructions (Hubble, Ricci, Granda--Oliveros, and---we
  emphasize---the apparent horizon) have a delta-function kernel, cluster like
  an ordinary $c_s^2\!=\!1$ fluid, and remain growth-informative. Light-cone
  constructions (the future event horizon and the particle horizon) have a
  one-dimensional oscillatory kernel; phase-averaging of the potential along
  the ray suppresses their density contrast as $\dDE/\dm\sim1/k^{3}$, so they
  carry no growth information beyond the background expansion. Smooth-nonlocal
  constructions are smoother still. We give the criterion, prove the partition
  is exhaustive, identify the future event horizon as the unique cutoff whose
  growth channel is \emph{unconditionally} closed, and delineate the two
  boundary subtleties---acoustic resonance on the particle-horizon ray, and a
  crossover scale for cutoffs with a small local admixture---that separate the
  exact statement from its observational version. The result reorganizes a large
  and fragmented model space around a single computable diagnostic and tells an
  observer, in advance, which geometric dark-energy models a growth measurement
  can and cannot constrain.
  \end{abstract}
  \vspace{1em}
  \end{@twocolumnfalse}
]

%========================================================================
\section{Introduction}

Dark energy is increasingly probed not only through the background expansion
$H(z)$ but through the rate at which cosmic structure grows, $f\sigma_8(z)$.
The premise of that program is that growth carries information the expansion
does not: two models with identical $H(z)$ can differ in how dark energy
clusters, and $f\sigma_8$ sees the difference. For a broad and popular class of
models, however, that premise fails---and it fails for a structural reason that
can be stated once and checked model by model.

The class in question is what we will call \emph{geometric dark energy}: models
in which the dark-energy density is not an independent dynamical field but is
fixed by a construction built from the spacetime geometry. Holographic dark
energy, $\rho_{\rm DE}\propto L^{-2}$ with $L$ an infrared cutoff length, is the
best-known example, and its generalized-entropy variants (Barrow, Tsallis,
R\'enyi, Kaniadakis) form a large post-DESI literature. But the category is
wider: it includes Ricci and Granda--Oliveros dark energy, apparent-horizon and
thermodynamic constructions, and nonlocal-gravity models in which the
dark-energy density is a nonlocal functional of the curvature.

For every such model the same fact holds: the dark-energy density perturbation
is not free. If $\rho_{\rm DE}=f(L[g])$ for some geometric length functional
$L$ of the metric $g$, then
\begin{equation}
\dDE=\frac{d\ln f}{d\ln L}\,\frac{\delta L}{L},
\label{eq:inherit}
\end{equation}
and the entire scale dependence of $\dDE$ lives in $\delta L$, the perturbation
of the geometric construction. The prefactor---where all the entropy physics
sits---is a bounded, scale-independent number. Whether growth can see the model
is therefore not a question about entropy or equation of state; it is a
question about $\delta L$, and specifically about \emph{how $\delta L$ responds
to the gravitational potential as a function of scale}.

This paper makes that question precise and answers it in general. We show
(Sec.~\ref{sec:kernel}) that $\delta L$ is always a linear functional of the
metric perturbation, so that in Fourier space
$\delta L_k=\int \mathcal{K}(k;\cdot)\,\Phi\,$, and that the observability of the
model is fixed by the large-$k$ scaling of the kernel $\mathcal{K}$. We prove
(Sec.~\ref{sec:partition}) that geometric dark-energy models partition
exhaustively into a small number of classes by this scaling, and that only the
local class remains growth-informative at observable scales. We identify
(Sec.~\ref{sec:eventhorizon}) the future event horizon as the unique
construction whose channel is unconditionally closed, and we treat
(Sec.~\ref{sec:boundaries}) the two boundary cases---particle-horizon acoustic
resonance and mixed cutoffs---that mark the edge of the clean statement. The
holographic closed-channel result established for event-horizon models
elsewhere~\cite{CompanionPaper} is the central, cleanest member of the
classification given here.

%========================================================================
\section{The response kernel}
\label{sec:kernel}

We work in the conformal-Newtonian gauge with metric potential $\Phi$, on
sub-horizon scales $k\gg aH$ where the observations live, and to linear order.
Matter perturbations obey the standard growth equation, whose only cosmological
input is $H(a)$; two models with the same expansion history give the same
matter growth unless dark energy itself clusters and back-reacts through the
Poisson equation,
\begin{equation}
-\frac{k^2}{a^2}\Phi=4\pi G\left[\rho_{\rm m}\dm+\rho_{\rm DE}\dDE(1+3c_{\rm eff}^2)\right].
\label{eq:poisson}
\end{equation}
The dark-energy term is the whole question. Its size relative to the matter
term is governed by $\dDE$, which by Eq.~\eqref{eq:inherit} is set by $\delta L$.

\paragraph{$\delta L$ is a linear functional of $\Phi$.}
A geometric length $L$ is a functional of the metric. Its first-order
perturbation is therefore a linear functional of the metric perturbation:
schematically
\begin{equation}
\delta L(x,t)=\int \mathcal{K}(x,t;x',t')\,\Phi(x',t')\,d^4x',
\label{eq:linfunctional}
\end{equation}
for a kernel $\mathcal{K}$ determined entirely by how $L$ is built from $g$.
In Fourier space $\delta L_k=\widetilde{\mathcal{K}}(k)\,\Phi_k$ (times a
possible integral along a curve or over a region, made explicit below).
Sub-horizon, the potential is sourced by matter, $\Phi_k\sim -4\pi G a^2
\rho_{\rm m}\dm/k^2$, so the density contrast carries the Poisson factor
$1/k^2$ together with the ultraviolet scaling of the kernel,
\begin{equation}
\frac{\dDE}{\dm}\;\sim\;\frac{1}{k^{2}}\;S(k),\qquad
S(k)\equiv\frac{\widetilde{\mathcal K}(k)}{\widetilde{\mathcal K}_{\rm local}},
\label{eq:master}
\end{equation}
with $S$ normalized to unity for a purely local response. The organizing claim
of this paper is that the growth-observability of a geometric dark-energy model
is fixed by $S(k)$---what the geometry adds beyond the universal pressure
(Jeans) floor $1/k^2$. Section~\ref{sec:partition} makes both factors precise:
the floor is the same for every construction (Lemma~\ref{lem:factor}), while
$S(k)$ takes one of a few forms set by the type of geometric functional
(Lemma~\ref{lem:kernels}). We caution that the local class does \emph{not}
follow from a naive derivative count on $\widetilde{\mathcal K}$; its value is
set by the pressure floor, as the proof makes explicit.

\paragraph{$S(k)$ is gauge-invariant.}
Although $\delta L$ for a nonlocal length is gauge-dependent---it integrates the
metric perturbation over a curve or region, and a coordinate change shifts the
integrand---the ultraviolet scaling of $S(k)$ is not. Two facts secure this.
First, the potential appearing in Eq.~\eqref{eq:linfunctional} is the
Newtonian-gauge $\Phi$, which equals the Bardeen invariant $\Phi_A$; the
geometric object over which it is integrated (a null ray, a causal region, the
support of $\Box^{-1}$) is coordinate-independent, so the leading large-$k$
behaviour of $\widetilde{\mathcal K}$ is a property of the invariant pair
$(\Phi_A,\text{geometry})$, not of the slicing. Second, $S$ is defined as a
\emph{ratio} to the local response, and any residual gauge contribution to
$\delta L$ enters at $\mathcal{O}(aH/k)$ relative to the leading term---the same
order in every gauge---so it cancels in the ratio to the order at which $S(k)$
is evaluated. A gauge transformation can therefore rescale $\delta L$ but cannot
change the exponent of $S(k)$; the local/nonlocal boundary of
Lemma~\ref{lem:kernels}, which depends only on that exponent, is a
gauge-invariant statement. We verified this explicitly for the light-cone class
by recomputing the leading scaling with $\Phi_A$ in place of the Newtonian
$\Phi$: the $1/k$ is unchanged.

%========================================================================
\section{The classification}
\label{sec:partition}

The kernel $\widetilde{\mathcal K}(k)$ is fixed by the \emph{type} of geometric
object $L$ is. There are only a few types, distinguished by the support of
$\mathcal K$ in Eq.~\eqref{eq:linfunctional}, and they exhaust the possibilities
because every geometric $L$ is one of them.

\begin{definition}[Kernel classes]
(Table~\ref{tab:classes} places the standard constructions in these classes.)
A geometric length functional $L$ is
\begin{itemize}
\item[\textbf{(L)}] \emph{local} if $\mathcal K$ is a delta function (with
possible derivatives): $\delta L=c(t)\Phi+\dots$. Then
$\widetilde{\mathcal K}\sim k^0$.
\item[\textbf{(N)}] \emph{null-line} if $\mathcal K$ has one-dimensional support
along a null ray: $\delta L_k=\int_{\rm ray}\Phi_k(u)\,e^{iku}\,du$. Then
$\widetilde{\mathcal K}\sim 1/k$ for a smooth envelope.
\item[\textbf{(S)}] \emph{smooth-nonlocal} if $\mathcal K$ is a rapidly
decaying spatial window (e.g. a Gaussian average over scale $R$,
$\widetilde{\mathcal K}\sim e^{-k^2R^2}$).
\item[\textbf{(I)}] \emph{inverse-box nonlocal} if $L$ involves
$\square^{-1}$ acting on curvature, so $\widetilde{\mathcal K}\sim 1/k^2$.
\end{itemize}
\end{definition}

\begin{theorem}[Exhaustive growth partition]
\label{thm:partition}
Let dark energy be geometric, $\rho_{\rm DE}=f(L[g])$ with $d\ln f/d\ln L$
bounded and nonzero. Then its sub-horizon density contrast relative to matter
scales as
\begin{equation}
\frac{\dDE}{\dm}\sim
\begin{cases}
1/k^{2} & \text{class (L): clusters, growth-informative}\\[2pt]
1/k^{3} & \text{class (N): growth-empty}\\[2pt]
\text{$\ll 1/k^{n}\,\forall n$} & \text{class (S): super-smooth}\\[2pt]
1/k^{4} & \text{class (I): growth-empty}
\end{cases}
\label{eq:partition}
\end{equation}
and every geometric $L$ is growth-informative or growth-empty according to the
dichotomy of Lemma~\ref{lem:kernels}; the four rates are those of the
constructions in current use.
\end{theorem}

The classification rests on a decomposition of the density response into a
universal part and a geometry-dependent part.

\begin{lemma}[Jeans floor and suppression factor]
\label{lem:factor}
For geometric dark energy with $\rho_{\rm DE}=f(L[g])$, the sub-horizon
density contrast factorizes as
\begin{equation}
\frac{\dDE}{\dm}\;\sim\;\underbrace{\frac{1}{k^{2}}}_{\text{Jeans floor}}\;\times\;S(k),
\label{eq:factor}
\end{equation}
where $S(k)$ is the ultraviolet scaling of the kernel $\widetilde{\mathcal
K}(k)$ through which $\delta L$ responds to the potential, normalized so that a
purely local response gives $S=1$.
\end{lemma}

\begin{proof}
Two facts combine. \emph{(i) The floor.} We take as the defining property of
\emph{geometric} dark energy that its density is a functional of the metric with
no independent matter or scalar degree of freedom---no separate kinetic term in
the action carrying its own momentum. For such a component the rest-frame sound
speed is not a free Lagrangian parameter; the prescription that renders
$\rho_{\rm DE}=f(L[g])$ well-defined ties its pressure perturbation to its
density perturbation through the same geometric relation, fixing the effective
sound speed to the smooth branch $c_{\rm eff}^2=1$. A component with
$c_{\rm eff}^2=1$ has its sound horizon equal to the Hubble radius, so
sub-horizon its perturbations are pressure-supported and Jeans-suppressed,
$\dDE/\dm\sim(aH/k)^2=1/k^2$ (in units $aH=1$)~\cite{Hu1998,WellerLewis2003}.
This is the universal floor for the geometric class, independent of which
construction is used. We flag the one way out explicitly: a functional that
implicitly reintroduces an independent momentum flux---for instance by making
$L$ depend on the matter velocity field, not the metric alone---would carry its
own $c_{\rm eff}^2\neq1$ and fall outside this class; our floor, and hence the
partition, is asserted for metric functionals, which is the case of interest and
covers every construction in Table~\ref{tab:classes}. \emph{(ii) The factor.} By
Eq.~\eqref{eq:inherit}, $\dDE=(d\ln f/d\ln L)\,\delta L/L$ with the prefactor
bounded, and by Eq.~\eqref{eq:linfunctional} $\delta L_k=\widetilde{\mathcal
K}(k)\,\Phi_k$. Any suppression of $\widetilde{\mathcal K}$ beyond the local
(delta-function) response multiplies the floor: writing $S(k)\equiv
\widetilde{\mathcal K}(k)/\widetilde{\mathcal K}_{\rm local}$, the response is
the floor times $S(k)$. A local response has $S=1$ and saturates the floor; a
nonlocal one has $S(k)\to0$ and lies below it. The two statements are
independent---(i) is a fluid pressure result, (ii) a functional-analytic
scaling---so their product is Eq.~\eqref{eq:factor}.
\end{proof}

\begin{lemma}[Locality dichotomy]
\label{lem:kernels}
On the homogeneous background the response $\Phi\mapsto\delta L$ is linear and
translation-invariant, hence diagonal in Fourier space,
$\delta L_k=\widetilde{\mathcal K}(k)\,\Phi_k$, and the suppression factor
$S(k)=\widetilde{\mathcal K}(k)/\widetilde{\mathcal K}_{\rm local}$ has a
well-defined ultraviolet order. Then:
\begin{itemize}
\item[(i)] If $\mathcal K$ contains a local (delta-supported) component---$L$
depending on the metric and finitely many derivatives at the point---that
component dominates the ultraviolet, $S(k)\!\to\!\mathrm{const}$ or grows, so
$\dDE/\dm\gtrsim1/k^{2}$: the dark energy clusters at or above the Jeans floor
and remains growth-informative.
\item[(ii)] If $\mathcal K$ is genuinely nonlocal---an integrable kernel with no
delta component---then by the Riemann--Lebesgue lemma
$\widetilde{\mathcal K}(k)\!\to\!0$, so $S(k)\!\to\!0$ and
$\dDE/\dm=o(1/k^{2})$: the dark energy is smoother than the Jeans floor and is
growth-empty.
\end{itemize}
The boundary $S\sim k^0$ is therefore exactly the local/nonlocal divide.
\end{lemma}

\begin{proof}
Linearity and translation invariance on an FRW background give the Fourier
diagonalization directly. For (i), a delta-supported kernel has
$\widetilde{\mathcal K}(k)=\sum_n c_n(ik)^n$, which does not vanish at large $k$;
the leading term sets $S$, and the density response is bounded below by the
$c_{\rm eff}^2=1$ Jeans floor of Lemma~\ref{lem:factor}(i), since a local
$\dDE\sim\Phi$ is a genuine pressure-supported fluid perturbation (the extra
derivative terms of a curvature cutoff enter only at subleading order
$(aH/k)^2$ and do not lift the response above the floor's scaling). For (ii),
an $L^1$ kernel with no atomic part has $\widetilde{\mathcal K}(k)\to0$ by
Riemann--Lebesgue, hence $S(k)\to0$; the model is strictly smoother than a
local one at every large $k$.
\end{proof}

The dichotomy is the rigorous content: \emph{growth can see a geometric
dark-energy model if and only if its density responds to the potential through
a locally-supported kernel.} Intuitively: a local response reads the potential
at a point, so a small-scale fluctuation registers at full strength; a nonlocal
response averages the potential over an extended ray or region, and a
small-scale fluctuation---many crests and troughs within the averaging
domain---largely cancels itself out. The more spread-out the kernel, the more
complete the cancellation, which is why the smooth spatial window (class S) is
smoother than the light-cone integral (class N), which is smoother than the
local cutoff (class L). Riemann--Lebesgue is the precise form of ``averaging an
oscillation sends it to zero.'' The specific exponents of
Theorem~\ref{thm:partition} refine this by computing the \emph{rate} at which
$S(k)$ vanishes for the nonlocal constructions actually in use---a rate that
Riemann--Lebesgue alone does not fix. For a light-cone integral, integration by
parts on the oscillatory kernel gives $S\sim1/k$ (a smooth envelope;
Sec.~\ref{sec:boundaries}), hence $1/k^3$. For an inverse d'Alembertian,
$S\sim1/k^2$, hence $1/k^4$. For a smooth spatial window, $S$ decays faster than
any power (Paley--Wiener). We do not claim these exhaust all conceivable
covariant functionals; they are the constructions the literature employs, and
each is placed by direct computation of its kernel order in
Table~\ref{tab:classes}.

\begin{proof}[Proof of Theorem~\ref{thm:partition}]
Lemma~\ref{lem:factor} gives $\dDE/\dm\sim S(k)/k^2$, and
Lemma~\ref{lem:kernels} establishes the dichotomy: $S\!\to\!\mathrm{const}$ for
a locally-supported kernel (clustering, growth-informative) and $S\!\to\!0$ for
a nonlocal one (growth-empty). The four entries of
Eq.~\eqref{eq:partition} are the values of $S(k)/k^2$ for the constructions in
Table~\ref{tab:classes}, using the kernel rates computed after
Lemma~\ref{lem:kernels}: $S=1$ (local) gives $1/k^2$; $S=1/k$ (light-cone) gives
$1/k^3$; $S=1/k^2$ (inverse-box) gives $1/k^4$; a smooth window gives faster
than any power. The partition into growth-informative and growth-empty is
exhaustive by the dichotomy; the specific exponents catalog the constructions
in use rather than claiming to exhaust all covariant functionals.
\end{proof}

Table~\ref{tab:classes} places the standard constructions. The entry that most
often causes confusion is the apparent horizon: in a spatially flat
Friedmann--Robertson--Walker background the apparent horizon coincides with the
Hubble radius, $L=1/H$, which is a \emph{local} function of the expansion rate.
Apparent-horizon dark energy is therefore in class (L)---it clusters---despite
the ``horizon'' name that might suggest a light-cone construction. Only the
future event horizon and the particle horizon, defined by null-geodesic
integrals, are class (N).

\begin{table}[!t]
\caption{Geometric dark-energy constructions by kernel class. Only class (L)
is informative to linear growth at observable scales.}
\label{tab:classes}
\begin{center}\small
\begin{tabular}{lll}\hline\hline
Construction & $L\sim$ & Class \\ \hline
Hubble cutoff & $1/H$ & (L) \\
Ricci cutoff & $(\dot H+2H^2)^{-1/2}$ & (L) \\
Granda--Oliveros & $(\alpha H^2+\beta\dot H)^{-1/2}$ & (L) \\
Apparent horizon (flat) & $1/H$ & (L) \\
Future event horizon & $a\!\int_t^\infty\!dt'/a$ & (N) \\
Particle horizon & $a\!\int_0^t\!dt'/a$ & (N) \\
Gaussian-averaged scale & smoothing window & (S) \\
Nonlocal $\square^{-1}R$ & inverse d'Alembertian & (I) \\
\hline\hline
\end{tabular}
\end{center}
\end{table}

%========================================================================
\section{The event horizon is unconditionally closed}
\label{sec:eventhorizon}

Among the growth-empty class (N), the future event horizon is the cleanest
case. Its defining ray runs from the observation time into the infinite future,
a region that (for accelerating cosmologies) is smoothly dark-energy dominated:
the potential envelope along the ray is monotone and non-oscillatory, so the
phase-averaging argument applies with no envelope condition at all, and
$\dDE/\dm\sim1/k^3$ holds for every event-horizon construction---every
generalized entropy, every deformation parameter---because the entropy enters
only the bounded prefactor in Eq.~\eqref{eq:inherit}. The particle horizon
reaches the same $1/k^3$ but by a slightly stronger route: its ray does cross
the oscillatory radiation era, yet the phase-average survives because the ray
is null and cannot resonate with the sub-luminal acoustic oscillation
(Sec.~\ref{sec:boundaries}). Class (N) is thus closed throughout, with the
event horizon closed \emph{unconditionally} and the particle horizon closed
by the light-versus-sound separation. This is the sense in which the
holographic closed-channel result~\cite{CompanionPaper} is not one result among
many but the representative of a class.

It is worth noting that this same cutoff is the one the data independently
prefer. Combined expansion-plus-growth analyses disfavor the Ricci and
Granda--Oliveros (class (L)) constructions because they overproduce early dark
energy, while the future-event-horizon construction remains
viable~\cite{Akhlaghi2018}. The cutoff for which the growth channel is provably
empty is thus also the cutoff that survives the data---so the closure is not a
statement about a disfavored corner of model space but about the observationally
central one.

%========================================================================
\section{Two boundary subtleties}
\label{sec:boundaries}

The clean partition of Theorem~\ref{thm:partition} has two edges where the exact
mathematical statement and its observational version diverge. Both are benign,
and naming them precisely is part of the result.

\paragraph{Acoustic resonance on the particle-horizon ray.}
The particle horizon is class (N), but its ray threads the radiation era, where
the potential is not a smooth envelope but oscillates with the acoustic phase.
Using the exact radiation-era transfer function
$\Phi\propto[\sin x-x\cos x]/x^3$ with $x=c_s k\tau$, the large-$x$ envelope
oscillates at frequency $c_s k$, while the null kernel oscillates at $\gamma k$,
where $\gamma$ is the ray's propagation speed in the relevant coordinate. Their
product beats to zero frequency---defeating the phase-average and degrading the
suppression from $1/k^3$ to $1/k^2$---only if
\begin{equation}
\gamma=c_s,
\end{equation}
that is, only if the null ray propagates at the acoustic sound speed. It does
not: the ray is null, $\gamma$ is the speed of light, and light is strictly
faster than sound in the photon-baryon plasma, $c_s=1/\sqrt3<1$. The
resonance is therefore not a fine-tuned near-miss but a structural
impossibility, and a high-precision evaluation of the ray integral with the full
transfer function confirms $\widetilde{\mathcal K}\sim1/k$ (hence
$\dDE/\dm\sim1/k^3$), robustly and independently of the ray parametrization.
The particle-horizon channel is thus closed for the same structural reason as
the event-horizon channel: the only configuration that would reopen it requires
sound to overtake light.

\paragraph{Envelope regularity along the ray.}
The class-(N) scaling $\widetilde{\mathcal K}\sim1/k$ was stated for a smooth
envelope, which invites the objection that a sharp background feature---a rapid
phase transition, a sudden shift in the dominant component---makes $\Phi(u)$
along the ray non-smooth and could, through the boundary terms of the
integration by parts, degrade the suppression back toward class (L). It cannot.
Integration by parts on $\int\Phi(u)e^{iku}du$ converts each order of
smoothness into one power of $1/k$: a jump discontinuity in the envelope
contributes a surface term of order $1/k$ (oscillating in $k$ through a phase
$e^{iku_0}$), a kink (derivative discontinuity) contributes $1/k^2$, and a
smooth envelope gives the interior $1/k$. In every case $\widetilde{\mathcal
K}=\mathcal{O}(1/k)$ or smaller. A direct evaluation with an imposed jump and an
imposed kink confirms $\widetilde{\mathcal K}\sim k^{-1}$ unchanged. Lifting the
model to class (L) ($S\sim k^0$) would require an $\mathcal{O}(1)$, non-decaying
$\widetilde{\mathcal K}$---that is, a delta function in $\Phi$ itself, a singular
potential, which is unphysical for any finite-energy perturbation. The honest
statement is therefore that the smooth-envelope assumption can be relaxed to
$\Phi$ of bounded variation along the ray, under which $\widetilde{\mathcal
K}=\mathcal{O}(1/k)$ and the class-(N) suppression is at worst $1/k^3$; sharp
features can only make it faster, never slower, and never local.

\paragraph{Mixed cutoffs and the crossover scale.}
Class membership is a statement about the \emph{pure} construction. A cutoff
with a local admixture, $L=\alpha L_{\rm EH}+\beta/H$, has
$\delta L_k=\alpha\,O(1/k)+\beta\,O(1)$: the local piece carries no
$k$-suppression and dominates once $k>k_{\rm cross}\equiv\alpha/\beta$. Any
nonzero $\beta$ therefore returns the model to class (L) asymptotically---but
only beyond $k_{\rm cross}$. For a small admixture ($\beta\ll\alpha$) the
crossover lies far above the observable range $k/aH\sim1$--$10^3$, and the model
is \emph{effectively} growth-empty across the window a survey can measure. The
exact boundary ($\beta=0$, unconditional $1/k^3$) and the observational boundary
(``dominantly null,'' $k_{\rm cross}$ beyond reach) thus differ, in exactly the
way an exact degeneracy differs from a precision-limited one. We state both.

The comparable-weight case $\alpha\sim\beta$ does not weaken the result, for two
reasons. First, it is not a marginal case for the theorem: with $\beta\sim\alpha$
one has $k_{\rm cross}\sim1$, so the model is class (L) across essentially the
entire observable range---it is \emph{growth-informative}, and the theorem
classifies it correctly and unambiguously as such. There is no regime the
partition fails to cover; the crossover simply tells us \emph{which} class a
mixed cutoff occupies as a function of scale. Second, a comparable local
admixture is disfavored on independent grounds: the $\beta/H$ Hubble-cutoff
piece reintroduces exactly the early-dark-energy overproduction that combined
expansion-plus-growth data disfavor for the pure Hubble and Ricci cutoffs
(Ref.~\cite{Akhlaghi2018}), scaled by $\beta$. A model with $\beta\sim\alpha$
inherits an $O(1)$ fraction of that disfavored behaviour and is constrained by
the same data that already disfavor class (L) horizon cutoffs. So the mixed
cutoff is either dominantly null (and effectively growth-empty, our case) or
carries a substantial local piece (and is both growth-informative \emph{and}
background-disfavored)---never a loophole that is simultaneously growth-hidden
and observationally viable.

%========================================================================
\section{Discussion}
\label{sec:discussion}

The classification converts a fragmented question---\emph{does this particular
horizon dark energy cluster?}---into a single computable diagnostic: identify
the kernel through which the density responds to the potential, read off its
ultraviolet scaling, and place the model in Eq.~\eqref{eq:partition}. The answer
requires no perturbation code and no fit; it follows from the form of the
geometric construction. For an observer, the practical content is a
prediction made in advance: a growth survey can constrain class (L) geometric
dark energy and cannot constrain classes (N), (S), or (I) beyond what the
background already fixes. For the holographic program specifically, it means the
event-horizon models that dominate the literature are, as a class, invisible to
$f\sigma_8$---and that the search for a growth signature should be redirected to
the local-cutoff constructions, where a signal can in principle exist.

Two limitations bound the scope. First, the framework is linear and
sub-horizon; the near-horizon and mildly nonlinear regimes, where the clean
$k$-scalings blur, are left for future work, though the observational weight of
$f\sigma_8$ sits at sub-horizon scales where the classification is sharp. We
note that nonlinear mode-coupling transfers power between scales and could in
principle feed the suppressed dark-energy contrast; but the coupling is through
the same potential $\Phi$, whose response is already classified, and the leaked
amplitude is bounded by the linear suppression times the (small) nonlinear
transfer efficiency---so a class-(N) model cannot acquire an $O(1)$ clustering
signal from nonlinearity, only a higher-order correction to an already-tiny
one. A quantitative nonlinear treatment is nonetheless the natural next check.
Second, the classification is of \emph{geometric} dark energy---models whose
density is inherited from the metric. It says nothing about scalar-field dark
energy (quintessence, k-essence), which carries a genuine Lagrangian sound speed
and an independent perturbation sector; such models are correctly outside the
partition, and the boundary between ``geometric'' and ``field-theoretic'' dark
energy is itself the natural next question.

Two further caveats deserve explicit mention. The class-(I) inverse-d'Alembertian
entry is placed by the leading Fourier scaling $\square^{-1}\sim1/k^2$ of the
nonlocal operator on the background; genuine nonlocal-gravity theories carry
their own causal boundary conditions and a richer operator structure, and we
classify only the leading UV scaling of the density response, not the full
dynamics of such theories---the $1/k^4$ entry should be read as ``at least as
smooth as class (N),'' with the precise treatment of a specific nonlocal-gravity
action left to a dedicated analysis. And the $c_{\rm eff}^2=1$ Jeans floor
assumes adiabatic initial conditions; a large dark-energy isocurvature mode
would alter the initial density scaling before the geometric kernel acts. For
geometric dark energy this is not expected---the density is slaved to the metric
and inherits the adiabatic mode---but a model with a genuinely independent
isocurvature initial condition would sit outside the floor's assumptions, worth
noting for completeness.

\begin{thebibliography}{9}
\bibitem{CompanionPaper} E. Schultz, ``Which generalized entropy? Discrimination
among horizon-entropy dark-energy models with DESI DR2 and Planck, and the
closure of the structure-growth channel,'' companion paper (2026).
\bibitem{Hu1998} W. Hu, ``Structure formation with generalized dark matter,''
Astrophys. J. \textbf{506}, 485 (1998), arXiv:astro-ph/9801234.
\bibitem{WellerLewis2003} J. Weller and A. M. Lewis, ``Large-scale cosmic
microwave background anisotropies and dark energy,'' Mon. Not. R. Astron.
Soc. \textbf{346}, 987 (2003), arXiv:astro-ph/0307104.
\bibitem{Akhlaghi2018} I. A. Akhlaghi, M. Malekjani, S. Basilakos, and H. Haghi,
``Model selection and constraints from holographic dark energy scenarios,''
Mon. Not. R. Astron. Soc. \textbf{477}, 3659 (2018), arXiv:1804.02989.
\end{thebibliography}

\end{document}
